\addcontentsline{toc}{subsubsection}{Solution 36. The Multiplicative Mirror: Positive Recurrence}

\begin{dlsummary}[title={\bfseries Solution 36. The Multiplicative Mirror: Positive Recurrence},breakable]

\begin{proof}
We prove that the only solutions are
\[
f(x)=x\qquad\text{and}\qquad f(x)=\frac1x.
\]

\medskip
\noindent\textbf{1. A sharp growth bound.}

Fix $t>1$ and write $a_n=f(t^n)$ for integers $n$. Substituting
$x=t$ and $y=t^n$ gives
\[
a_{n+1}+a_{n-1}\le (t+t^{-1})a_n. \tag{1}
\]
Put
\[
U_n=\frac{t^n-t^{-n}}{t-t^{-1}}.
\]
Iterating (1), and using positivity of all the terms, gives
\[
0<a_n\le U_n a_1-U_{n-1}a_0.
\]
Put
\[
b_n=a_{n+1}-\frac{a_n}{t}.
\]
The recurrence (1) rearranges to
\[
b_n\le t b_{n-1},
\]
so $b_n\le t^n b_0$. If $b_0<0$, then $b_n\to-\infty$, and
\[
a_{n+1}=\frac{a_n}{t}+b_n<0
\]
for all sufficiently large $n$, contradicting positivity. Hence
\[
a_1\ge\frac{a_0}{t}.
\]
Applying the same argument to the sequence $a_{-n}$ yields
\[
a_1\ge \frac{a_0}{t},\qquad a_{-1}\ge \frac{a_0}{t}.
\]
Since $a_0=f(1)=1$, (1) with $n=0$ also gives
$a_1+a_{-1}\le t+t^{-1}$. Consequently $a_1,a_{-1}\le t$. As $t>1$ was arbitrary,
\[
\min\{u,u^{-1}\}\le f(u)\le\max\{u,u^{-1}\}
\qquad(u>0). \tag{2}
\]

\medskip
\noindent\textbf{2. The pointwise alternatives.}

Taking logarithms in (2) gives
\[
|\log f(u)|\le |\log u|.
\]
Applying this inequality to $f(u)$ and using $f(f(u))=u$ gives the reverse
inequality. Thus
\[
|\log f(u)|=|\log u|,
\]
so
\[
f(u)\in\{u,u^{-1}\}\qquad(u>0). \tag{3}
\]

\medskip
\noindent\textbf{3. The choice cannot switch.}

For $u>1$, let
\[
L=\{u>1:f(u)=u^{-1}\}.
\]
If $u\in L$, substitute $y=u$ and $x=z/u$ in the original inequality. Using
(2) gives
\[
f(z)\le
\frac{z/u+u/z}{u}
-\min\left\{\frac{u^2}{z},\frac{z}{u^2}\right\}. \tag{4}
\]
As $z\to u$, the right-hand side tends to $u^{-1}<u$. Hence, for $z$
sufficiently close to $u$, the alternative $f(z)=z$ in (3) is impossible;
therefore $L$ is open in $(1,\infty)$.

Conversely, if $u_k\in L$ and $u_k\to u>1$, use (4) with centre $u_k$
and $z=u$, and pass to the limit. This gives $f(u)\le u^{-1}$, so (3)
implies $u\in L$. Thus $L$ is closed in $(1,\infty)$.

Since $(1,\infty)$ is connected, $L$ is either empty or all of $(1,\infty)$.
Finally, the involution makes $f$ injective, which forces the corresponding
choice on $(0,1)$ as well. The two possibilities are therefore precisely
\[
f(x)=x\qquad\text{or}\qquad f(x)=\frac1x.
\]
Direct substitution shows that both functions satisfy the inequality with
equality.
\end{proof}
\end{dlsummary}

\begin{takeaways}
\begin{itemize}
\item Multiplicative orbits turn the two-variable inequality into a recurrence.
Positivity makes the recurrence rigid, and applying the involution twice
turns a one-sided estimate into an equality. The final connectedness argument
is needed to exclude mixed pointwise choices.
\end{itemize}
\end{takeaways}